% La figura del vídeo D9: una función continua que cruza el eje entre a y b.
% Se compila a PDF con `pdflatex bolzano.tex`; el arnés copia el PDF a la
% carpeta figures/ de «El teorema de Bolzano».
\documentclass[tikz,border=6pt]{standalone}
\begin{document}
\begin{tikzpicture}[scale=1.4, line cap=round]
  \draw[->, gray!70] (-0.3,0) -- (4.4,0) node[right, black] {$x$};
  \draw[->, gray!70] (0,-1.5) -- (0,2.1) node[above, black] {$y$};
  \draw[very thick, green!45!black, domain=0.5:3.8, samples=80]
    plot (\x, {0.2*(\x-2)^3 + 0.25*(\x-2)});
  \draw[dashed, gray] (0.5,0) -- (0.5,-1.05) (3.8,0) -- (3.8,1.62);
  \fill (0.5,-1.05) circle (1.8pt) (3.8,1.62) circle (1.8pt);
  \node[below] at (0.5,0) {$a$};
  \node[below] at (3.8,-0.05) {$b$};
  \fill[red!70!black] (2,0) circle (2.4pt);
  \node[above left, red!70!black] at (2,0) {$c$};
  \node[left] at (0.5,-1.05) {$f(a) < 0$};
  \node[right] at (3.8,1.62) {$f(b) > 0$};
\end{tikzpicture}
\end{document}
